\begin{minipage}[t]{0.535\linewidth}
{\sffamily\bfseries\small CÓMO USAR ESTE CUADERNO}\par\vspace{2pt}
\begin{enumerate}[leftmargin=1.35em,itemsep=1.2pt,topsep=0pt]
  \item Copia la \textbf{plantilla base} en un archivo \codigo{main.cpp}.
  \item Busca la receta según el problema en el \textbf{mapa rápido}.
  \item Pega el \texttt{struct} o las funciones de la receta.
  \item En \codigo{main}, lee los datos, llama a la función y muestra el
        resultado con \codigo{cout}.
\end{enumerate}
\vspace{2pt}
{\footnotesize
No hace falta entender todo el código: las recetas ya están probadas.
Cambia solamente lo que dice \textbf{Entrada} y \textbf{Salida}.
Cada receta incluye un \codigo{main} de ejemplo; ese \codigo{main} se puede
borrar cuando ya tengas el tuyo.\par}
\vspace{2pt}
{\sffamily\bfseries\small REGLA DE ORO}\par\vspace{1pt}
{\footnotesize
Si el problema pide posiciones y usa 1-based, imprime \codigo{pos + 1}.
Todos los índices de este cuaderno empiezan en 0.\par}
\vspace{1pt}
{\footnotesize\textbf{Si algo falla:} revisa la lectura, el 0/1-based,
el uso de \codigo{long long} y que el SAM reciba minúsculas.\par}
\end{minipage}\hfill
\begin{minipage}[t]{0.44\linewidth}
{\sffamily\bfseries\small MAPA RÁPIDO: ¿QUÉ NECESITO?}\par\vspace{2pt}
\begin{itemize}[leftmargin=1.1em,itemsep=0.8pt,topsep=0pt]
  \item ¿Aparece un patrón? \hfill \textbf{KMP}
  \item ¿Prefijos que también son sufijos? \hfill \textbf{KMP}
  \item ¿Patrón en texto, versión corta? \hfill \textbf{Z}
  \item ¿Palíndromos? \hfill \textbf{Manacher}
  \item ¿Comparar substrings rápido? \hfill \textbf{Hash}
  \item ¿Substrings distintas, ocurrencias, LCS? \hfill \textbf{SAM}
\end{itemize}
\vspace{2pt}
{\sffamily\bfseries\small COMPILAR Y EJECUTAR}\par\vspace{1pt}
{\ttfamily\fontsize{6.7}{8.2}\selectfont
g++ -O2 -std=c++17 prog.cpp -o prog\par
./prog < entrada.txt > salida.txt\par}
\vspace{1pt}
{\footnotesize La primera línea compila; la segunda ejecuta. Si no usas
archivos, ejecuta solamente \codigo{./prog}.\par}
\end{minipage}
